\documentclass{article}

\PassOptionsToPackage{numbers, sort, compress}{natbib}
\usepackage[preprint]{neurips_2025}

\usepackage[T1]{fontenc}
\usepackage[utf8]{inputenc}
\usepackage[english,russian]{babel}
\usepackage{tempora}
\usepackage{hyperref}
\usepackage{url}
\usepackage{booktabs}
\usepackage{nicefrac}
\usepackage{microtype}
\usepackage{xcolor}

%%%

\usepackage{subcaption}
\usepackage{graphicx}
\usepackage{multirow}
\usepackage{amsmath,amssymb,amsfonts}
\usepackage{amsthm}
\usepackage{mathrsfs}
\usepackage{xcolor}
\usepackage{textcomp}
\usepackage{manyfoot}
\usepackage{algorithm}
\usepackage{algorithmicx}
\usepackage{algpseudocode}
\usepackage{listings}

\newtheorem{theorem}{Theorem} % continuous numbers
%%\newtheorem{theorem}{Theorem}[section] % sectionwise numbers
%% optional argument [theorem] produces theorem numbering sequence instead of independent numbers for Proposition
\newtheorem{proposition}[theorem]{Proposition}%
\newtheorem{lemma}{Lemma}%
%%\newtheorem{proposition}{Proposition} % to get separate numbers for theorem and proposition etc.

\newtheorem{example}{Example}
\newtheorem{remark}{Remark}

\newtheorem{definition}{Definition}
\newtheorem{assumption}{Assumption}

%%%

\input{math_commands.tex}


\title{Предсказание гроккинга за счёт вычисления размерности аттрактора нейросети}


% The \author macro works with any number of authors. There are two commands
% used to separate the names and addresses of multiple authors: \And and \AND.
%
% Using \And between authors leaves it to LaTeX to determine where to break the
% lines. Using \AND forces a line break at that point. So, if LaTeX puts 3 of 4
% authors names on the first line, and the last on the second line, try using
% \AND instead of \And before the third author name.


\author{
  Nickolay Karlov\\
  MIPT\\
  Moscow, Russia\\
  \texttt{karlov.na@phystech.edu}\\
  \And
  Alexey Kravatskiy\\
  MIRIAI\\
  Moscow, Russia\\
  \texttt{kravatskii.a@miriai.org}\\
}


\begin{document}


\maketitle

\begin{abstract}
    Обучение глубоких моделей сопровождается генерацией множества взаимосвязанных временных рядов (loss, accuracy, weight parameters, learning rate). В данной работе выдвигается гипотеза о том, что процесс оптимизации нейронной сети может быть рассмотрен как вынужденная динамическая система обладающая аттрактором, соответствующим локальному минимуму функции потерь, а феномен отложенной генерализации (гроккинг) представляет собой структурный фазовый переход на низкоразмерный аттрактор. Предложен статистический метод, использующий оценку размерности Левиной-Бикеля (MLE Intrinsic Dimension), способный осуществлять раннее предсказание гроккинга и оценивать истинную алгебраическую сложность выученного алгоритма через отслеживание коллапса размерности 1-D логов обучения. Метод был проверен на задачах модульного сложения и композиции в группе перестановок $S_5$ с использованием архитектуры Omnigrok. Результаты демонстрируют, что гроккинг сопровождается устойчивым снижением размерности аттрактора, что подтверждает гипотезу о топологическом коллапсе как механизме отложенной генерализации.
\end{abstract}

\textbf{Keywords:} гроккинг, метод Левиной-Бикеля, фазовое пространство, динамические системы, оптимизация нейросетей.

\section{Введение}\label{sec:intro}


В классической статистической теории обучения центральной парадигмой является баланс между смещением и дисперсией (bias-variance tradeoff). Долгое время считалось, что если в процессе оптимизации эмпирический риск (ошибка на обучающей выборке) падает до нуля, а зазор обобщения остается большим или начинает расти, модель безвозвратно переобучается, что требует остановки обучения. Однако в 2022 году был открыт феномен гроккинга \cite{power2022grokking} — эффекта отложенной генерализации, при котором модель, находясь в состоянии переобучения на протяжении тысяч итераций, совершает фазовый переход и достигает идеальной точности на тестовой выборке. Этот феномен ставит под сомнение классические подходы к мониторингу обучения и порождает необходимость в новых методах раннего предсказания обобщающей способности.

В современной литературе сформировалось несколько подходов к изучению гроккинга. Первое направление, основанное на механистической интерпретируемости (mechanistic interpretability) \cite{nanda2023progress, chughtai2023toy}, исследует гроккинг как процесс формирования структурированных алгебраических представлений (например, Фурье-базисов) внутри матриц весов. Хотя этот «white-box» подход дает глубокое понимание механики, он вычислительно немасштабируем и требует создания уникальных метрик (progress measures) под каждую новую задачу. Второе направление фокусируется на макроскопической динамике ландшафта функции потерь («black-box» подход). В работах \cite{liu2022omnigrok, thilak2022slingshot} анализируется влияние нормы весов и эффекта «рогатки» (slingshot effect), а авторы \cite{notsawo2023predicting} предлагают использовать спектральные сигнатуры лосса (параметры Йорта) для предсказания феномена.

Однако существующие макроскопические методы обладают существенным ограничением: они фиксируют лишь статистические изменения параметров системы, но не отражают изменения \textit{количества степеней свободы} (внутренней сложности) самой модели в процессе обучения.


В данной работе предлагается новый статистический метод, использующий методы поиска размерности аттрактора для предсказания и описания гроккинга. Траектория оптимизации формализуется как система $\vw_{t+1} = F(\vw_t, \xi_t)$, где $\vw_t$ — параметры модели, а $\xi_t$ — внешнее воздействие, обусловленное использованием стохастического градиента. Для восстановления свойств этой высокоразмерной системы по скалярным логам $y_t \in \R$ используется теорема Такенса о вложении \cite{takens1981detecting}. Чтобы учесть нестационарность и стохастичность градиентного спуска, теоретическая база расширяется за счет теорем Старка о delay embeddings для систем с детерминированным и стохастическим форсингом \cite{stark1999delay, stark2003delay}.

\textbf{Центральная гипотеза исследования} заключается в том, что гроккинг — это структурный топологический переход. В фазе меморизации оптимизатор хаотично блуждает в высокоразмерном пространстве параметров, пытаясь заучить тренировочную выборку. В фазе генерализации $L_2$-регуляризация (weight decay) заставляет систему сбросить лишние степени свободы и кристаллизовать компактный алгебраический алгоритм. Топологически это эквивалентно сжатию аттрактора в подмногообразие низкой размерности.

Для проверки предложенной гипотезы и анализа феномена гроккинга в работе проводится серия вычислительных экспериментов. В качестве основного математического инструмента применяется оценка внутренней размерности фазового пространства методом максимального правдоподобия Левиной-Бикеля (MLE Intrinsic Dimension) \cite{levina2004maximum}. Обучение проводится на алгебраических датасетах с использованием минималистичной архитектуры однослойного Трансформера, предложенной в исследовании Omnigrok \cite{liu2022omnigrok}. В первом эксперименте решается задача модульного сложения (коммутативная группа), во втором — задача операции композиции в группе перестановок $S_5$ (неабелева группа). Подобный выбор тестовых сред позволяет протестировать предложенный метод как на способность осуществлять раннее предсказание отложенной генерализации, так и на чувствительность к изменению скрытой алгебраической сложности задачи.











\section{Постановка задачи}\label{sec:problem}

\textbf{Формализация динамической системы и наблюдателя.}
Рассмотрим нейросетевую модель $\gM$, процесс обучения которой в течение $T$ итераций (шагов оптимизатора) можно описать как неавтономную стохастическую динамическую систему. Пусть $\vw_t \in \sW \subset \R^D$ — вектор параметров (весов) модели на шаге $t \in \{1, \dots, T\}$. Правило обновления параметров имеет вид:
\begin{equation}
    \vw_{t+1} = F(\vw_t, \xi_t),
\end{equation}
где $F$ — оператор оптимизации (например, SGD с импульсом или AdamW), а $\xi_t$ — внешнее воздействие (стохастический форсинг), включающее выбор мини-батча и расписание гиперпараметров.

На практике полное высокоразмерное состояние $\vw_t$ ($D \gg 1$) недоступно для вычислительно эффективного анализа. Наблюдаемой выборкой $\gD = \{ \vm_t \}_{t=1}^T$ является набор одномерных скалярных временных рядов $\vm_t =[m_1(t), \dots, m_K(t)]^\top$, генерируемых в процессе оптимизации. Каждая макроскопическая метрика (функция потерь, норма весов, норма градиентов) выступает в роли функции наблюдения $m_i(t) = \phi_i(\vw_t, \xi_t) \in \R$.



\textbf{Задача исследования:} Опираясь исключительно на одномерные ряды наблюдений $m_i(t)$, реконструировать топологические свойства аттрактора скрытой системы $\vw_t$ и разработать ранний предиктор фазового перехода оптимизатора.




\textbf{Определение феномена (Гроккинг).}
В качестве целевого фазового перехода рассматривается феномен отложенной генерализации (гроккинг). Пусть $\text{Acc}_{train}(\vw_t)$ и $\text{Acc}_{val}(\vw_t)$ — точность модели на обучающей и валидационной выборках. Введем два критических момента времени для заданного малого порога ошибки $\epsilon > 0$:
\begin{itemize}
    \item \textbf{Момент меморизации ($t_{mem}$):} шаг, на котором сеть впервые идеально интерполирует обучающую выборку, $\text{Acc}_{train}(\vw_{t_{mem}}) \ge 1 - \epsilon$, при этом валидационная точность остается на уровне случайного угадывания $\text{Acc}_{val}(\vw_{t_{mem}}) \approx 0$.
    \item \textbf{Момент генерализации ($t_{gen}$):} шаг, на котором происходит скачкообразный рост обобщающей способности, $\text{Acc}_{val}(\vw_{t_{gen}}) \ge 1 - \epsilon$.
\end{itemize}
Процесс классифицируется как гроккинг при условии аномальной задержки: $t_{gen} \gg t_{mem}$.

\textbf{Критерий качества предсказания.}
Основная гипотеза работы: переход от меморизации к генерализации является процессом структурного схлопывания степеней свободы оптимизатора. Пусть $\hat{d}(t)$ — эффективная размерность реконструированного аттрактора в момент времени $t$. Поскольку мгновенные оценки могут содержать вычислительный шум, успешным предсказанием (early warning signal) гроккинга считается формирование устойчивого нисходящего тренда предшествующего гроккингу и завершающегося стабилизацией размерности на новом низком уровне:
\begin{equation}
    \exists t_{drop} \in (t_{mem}, t_{gen}) \quad \text{такое, что} \quad \forall t \in [t_{drop}, t_{gen}]: \;\; \hat{d}(t) \ll \hat{d}(t_{mem}).
\end{equation}
Дополнительным условием является то, что на интервале $(t_{mem}, t_{drop})$ должно наблюдаться в среднем монотонное падение $\hat{d}(t)$ без значимых возвратных выбросов вверх. Иными словами, статистически значимый обвал эффективной размерности аттрактора и её фиксация на новом низкоразмерном плато должны произойти \textit{строго до} фактического роста валидационной точности.


% \textbf{Contributions.} Our contributions can be summarized as follows:
% \begin{itemize}
%     \item We present...
%     \item We demonstrate the validity of our theoretical results through empirical studies...
%     \item We highlight the implications of our findings for...
% \end{itemize}

% \textbf{Outline.} The rest of the paper is organized as follows...

% \section{Related Work}\label{sec:rw}

% \textbf{Topic \#1.}
% TODO

% \textbf{Topic \#2.}
% TODO

% \section{Preliminaries}\label{sec:prelim}

% \subsection{General notation}

% In this section, we introduce the general notation used in the rest of the paper and the basic assumptions.

% \subsection{Assumptions}

% TODO

\section{Теоретическое обоснование и предлагаемый метод}\label{sec:method}

Для решения поставленной задачи предложен метод, опирающийся на теорию нелинейных динамических систем. В данном разделе обосновывается возможность извлечения геометрии многомерной системы из 1D-лога и приводится алгоритм оценки эффективной размерности $\hat{d}(t)$.

\subsection{Обучение как вынужденная динамическая система (Теоремы о вложении)}

Согласно классической теореме Такенса \cite{takens1981detecting}, для автономных систем, сходящихся к аттрактору истинной размерности $m$, топологию всего аттрактора можно реконструировать по одному скалярному ряду наблюдений $X(t) \in \{m_i(t)\}$, используя векторы запаздывания (delay embeddings):
\begin{equation}
    \vx_t =[X(t), X(t-\tau), \dots, X(t-(E-1)\tau)]^\top \in \R^E,
\end{equation}
при условии, что размерность пространства вложения $E \ge 2m + 1$.

Фундаментальным математическим требованием теорем о вложении является \textit{«типичность»} (genericity) функции наблюдения $\phi(\cdot)$. С точки зрения дифференциальной топологии это означает, что функция метрики не должна обладать вырожденными симметриями, которые склеивают различные точки фазового пространства при проекции. Например, агрегированная норма весов $||\vw_t||_2$ обладает сферической симметрией, измеряя лишь радиальное расстояние от начала координат. В определенных режимах она может выступать в роли «вырожденного наблюдателя» (degenerate observer), теряющего информацию об алгебраической структуре (угловых координатах) весов. В то же время, нелинейные метрики, напрямую зависящие от распределения логитов (такие как функция потерь $\Ls$), выступают в роли типичных наблюдателей, сохраняющих многомерную топологию системы.

Однако градиентный спуск в глубоком обучении не является автономной системой. Влияние случайного выбора мини-батчей, а также использование расписаний гиперпараметров (learning rate schedulers) формально нарушают классические условия Такенса. Для математического обоснования нашего подхода мы опираемся на теоремы Старка о вложениях для вынужденных систем со смешанным воздействием \cite{stark1999delay, stark2003delay}.

В современной теории оптимизации процесс обучения $\vw_{t+1} = F(\vw_t, \omega_t, \xi_t)$ рассматривается под воздействием двух типов форсинга. Детерминированный форсинг $\omega_t$ (расписание шага обучения, искусственные вмешательства в батч) подчиняется теореме Старка для детерминированных систем \cite{stark1999delay}, гарантирующей сохранение топологии при гладком внешнем драйвере. В свою очередь, стохастический форсинг $\xi_t$ (градиентный шум мини-батчей) аппроксимируется многомерным Гауссовским процессом, а выбор батча из конечного набора описывается в терминах систем итерируемых функций (Iterated Function Systems) \cite{stark2003delay}. Для обоих случаев Старк доказал, что отображение задержки сохраняет свойства диффеоморфизма. Это фундаментально обосновывает правомерность измерения размерности скрытого аттрактора весов $\sW$, оперируя исключительно скалярными логами в реконструированном фазовом пространстве $\R^E$, независимо от наличия детерминированных расписаний или стохастического шума.

\subsection{Оценка эффективной размерности фазового пространства}\label{sec:theory_dim}

Для детекции фазового перехода нам необходимо оценивать размерность аттрактора системы. В нашем методе мы рассмотрели как классические, так и современные методы, ранжировав их по устойчивости к специфике нейросетевой оптимизации (шуму мини-батчей):

\textbf{1. Метод ложных ближайших соседей (FNN) \cite{kennel1992determining} и Метод Цао \cite{cao1997practical}:}
Геометрические подходы "снизу-вверх". FNN проверяет наличие ложных пересечений траектории при переходе от размерности $d$ к $d+1$, используя жесткие пороговые эвристики ($R_{tol}$). Метод Цао устраняет потребность в порогах, анализируя усредненное отношение расстояний. Однако оба метода требуют глобальной возвратности траектории (рекуррентности), из-за чего они демонстрируют низкую стабильность на транзиентных (переходных) участках оптимизации, часто вырождаясь в тривиальную оценку $\hat{d}=1$.

\textbf{2. Метод максимизации симплексной проекции \cite{sugihara1990nonlinear}:}
Оптимальное $E$ ищется как аргумент, максимизирующий авторегрессионную корреляцию $\rho$ при предсказании будущего состояния $y_{t+1}$ по его соседям в $\R^E$. Метод топологически гарантирует наилучшее развертывание аттрактора, однако его разрешение является целочисленным (дискретным), что загрубляет детекцию плавных фазовых переходов.

\textbf{3. Метод максимального правдоподобия (MLE Intrinsic Dimension) \cite{levina2004maximum}:}
Выбран в качестве основного аналитического инструмента в нашем методе. В отличие от геометрических методов, подход Левиной-Бикеля проецирует временной ряд сразу в избыточное высокоразмерное пространство (например, $E=15$) и оценивает локальную внутреннюю размерность $\hat{d}_{x}$ статистически, по Пуассоновскому распределению расстояний до $k$ ближайших соседей. Процесс формализован в Алгоритме \ref{alg:mle_id}.

\begin{algorithm}[htbp]
\caption{Оценка динамической размерности MLE (Levina-Bickel)}\label{alg:mle_id}
\begin{algorithmic}[1]
\Require Временной ряд наблюдений $\gD_{obs}$, задержка $\tau$, избыточная размерность $E_{max}$, окно $W$, соседи $k$
\Ensure Эффективная размерность $\hat{d}(t)$
\State $M \gets \text{DelayEmbedding}(\gD_{obs}, E_{max}, \tau)$ \Comment{Вложение в $\R^{E_{max}}$}
\For{каждой точки $\vx_t \in M$ внутри скользящего окна $W$}
    \State Найти $k$ ближайших соседей с расстояниями $r_1 \le r_2 \le \dots \le r_k$
    \State $\hat{d}_{x_t} \gets \left[ \frac{1}{k-1} \sum_{j=1}^{k-1} \log \left( \frac{r_k}{r_j} \right) \right]^{-1}$ \Comment{Локальная оценка по радиусу $r_k$}
\EndFor
\State $\hat{d}(t) \gets \frac{1}{|W|} \sum_{\vx_t \in W} \hat{d}_{x_t}$ \Comment{Усреднение размерности на окне $W$}
\State \Return $\hat{d}(t)$
\end{algorithmic}
\end{algorithm}

Данный метод возвращает \textit{непрерывную} величину, которая вычисляется локально, не требует глобальной возвратности траектории и демонстрирует устойчивость к шуму мини-батчей.

\subsection{Физика Гроккинга как топологического коллапса}\label{sec:theory_grokking}

Применяя описанный математический аппарат, мы формализуем гроккинг как структурный фазовый переход в ландшафте потерь.
\begin{itemize}
    \item \textbf{Фаза меморизации:} Сеть "заучивает" данные как несвязанный шум. Траектория оптимизатора $\vw_t$ блуждает по высокоразмерному многообразию. В этот период эффективная размерность $\hat{d}(t)$, вычисляемая алгоритмом MLE, будет максимальной.
    \item \textbf{Фаза генерализации (Гроккинг):} Под воздействием $L_2$-регуляризации (weight decay) сеть отбрасывает шумовые признаки и «кристаллизует» оптимальный алгебраический алгоритм в структуре своих весов.
\end{itemize}

Нахождение обобщающего решения строго ограничивает степени свободы системы. Геометрически это означает, что аттрактор динамической системы стягивается в низкоразмерное подмногообразие. Следовательно, отслеживание эффективной размерности $\hat{d}(t)$ на скользящем окне должно действовать как опережающий индикатор топологического упрощения системы.













\section{Вычислительный эксперимент}\label{sec:exp}

\subsection{Планирование эксперимента}\label{sec:exp_plan}

\textbf{Цели эксперимента.}Глобальной целью вычислительного эксперимента является проверка гипотезы о гроккинге как о структорном переходе сопровождающемся понижением размерности аттрактора и валидация статистического метода как раннего предиктора гроккинга. В исследование проводилось три эксперимента:
\begin{enumerate}
    \item \textbf{Базовый гроккинг на задаче модульного сложения:}
    В данном эксперименте мы используем обычные torch слои из-за чего гроккинг получается весьма не гладким.
    \item \textbf{Гроккинг архитектуры Omnigrok на задаче модульного сложения:}
    В данном эксперименте мы используем архитектуру Omnigrok, которая была специально разработана для демонстрации гроккинга. Она включает в себя кастомные слои и оптимизацию, что обеспечивает более плавный и отчетливый эффект гроккинга.
    \item \textbf{Гроккинг на задаче композиции в $S_5$:}
    В данном эксперименте мы также используем архитектуру Omnigrok. Задача композиции в группе перестановок $S_5$ является более сложной алгебраической задачей, что позволяет проверить чувствительность метода к изменению скрытой сложности задачи и его способность предсказывать гроккинг в более сложных условиях.
\end{enumerate}

\textbf{Наборы данных и конфигурация алгоритма.}
Для всех экспериментов использовался однослойный трансформер, для первого экперимента это был стандартный трансформер, а для второго и третьего эксперимента это была архитектура Omnigrok. В первом и втором экспериментах использовался датасет модульного сложения по модулям 97 и 113 соответственно, в третьем эксперименте использовался датасет композиции в $S_5$. Все модели обучались с помощью оптимизатора AdamW с фиксированным learning rate и weight decay.

Помимо базовых метрик (функция потерь $\Ls$, точность), согласно современным подходам к анализу обобщающей способности \cite{jiang2019fantastic} и нелинейной динамики \cite{kantz2004nonlinear}, алгоритм отслеживает дополнительные скалярные сигналы: норму градиента $\|g\|_2$, нормы весов $\|\vw\|_2$, а также скалярное произведение между градиентами на соседних итерациях. Гиперпараметры фазовой реконструкции (задержка $\tau$ и размерность вложения $E$) вычисляются динамически на основе минимума взаимной информации (DMI) \cite{fraser1986independent} и метода MLE соответственно.

\textbf{Ожидаемые графики и таблицы.}
Планируется получить график динамики размерности вложения $d^*(t)$ с явно выраженным спадом в момент, предшествующий резкому росту валидационной точности (гроккинг).

\subsection{Результаты экспериментов}\label{sec:exp_report}

\textbf{Результаты сценария Grokking.}
Исследование феномена отложенной генерализации проводилось в несколько этапов: от базового подтверждения гипотезы до тестирования метода в условиях сложной алгебраической структуры данных и стохастического шума.

\textbf{Этап 1: Базовый эксперимент и проблема вырожденного наблюдателя.}
Первоначальные эксперименты проводились на задаче модульного сложения с использованием полнопакетного градиентного спуска (Full-Batch GD) (Рис. \ref{fig:grokking_main}). Нам удалось зафиксировать падение эффективной размерности (MLE ID) ряда $\Ls_{train}$ непосредственно перед скачком валидационной точности (с $E \approx 4-5$ до $E \approx 2.8$).

Однако этот эксперимент выявил важную методологическую проблему. При полнопакетной оптимизации или замене Adam на SGD корреляция между размером фазового пространства $E$ и эффектом гроккинга становилась менее очевидной. С теоретической точки зрения это объясняется проблемой «вырожденного наблюдателя»: в отсутствие стохастического шума мини-батчей траектория оптимизатора вырождается в детерминированную гладкую кривую, а обучающий лосс падает практически до нуля, лишая метод MLE необходимой дисперсии для оценки истинной размерности пространтства $d^*(t)$.

\begin{figure}[H]
     \centering
     \begin{subfigure}[b]{0.42\textwidth}
         \centering
         \includegraphics[width=\textwidth]{images/grokking_accuracy.png}
         \caption{Динамика Train/Val Accuracy}
     \end{subfigure}
     \hfill
     \begin{subfigure}[b]{0.32\textheight}
         \centering
         \includegraphics[width=\textwidth]{images/grokking_dimension.png}
         \caption{Коллапс размерности (MLE ID)}
     \end{subfigure}
     \caption{Этап 1: Феномен гроккинга на базовом эксперименте. Переход от запоминания к обобщению сопровождается снижением топологической сложности траектории $\Ls_{train}$.}
     \label{fig:grokking_main}
\end{figure}

\textbf{Этап 2: Стохастическая оптимизация и не типичный наблюдатель (Модульное сложение).}
Для решения описанной проблемы мы перешли к стохастической оптимизации (Mini-batch AdamW) и расширили спектр анализируемых 1-D скалярных метрик. Стохастический шум (Stochastic Forcing) увеличил дисперсию рядов скалярных метрик, что позволило алгоритмам фазовой реконструкции более точно оценивать истинную размерность фазового пространства аттрактора. В качестве целевой метрики была выбрана норма весов ($||\vw||_2$), так как она сохраняет богатую дисперсию на протяжении всей фазы меморизации и что не менее важно, ни как не зависит от валидационной выборки.

На Рис. \ref{fig:grokking_flat} представлены результаты эксперимента. В верхней строке (с включенной $L_2$-регуляризацией, $WD=1.0$) наблюдается идеальный топологический коллапс: эффективная размерность $E$(а точнее её оценка $d^*(t)$ полученная MLE методом) обваливается синхронно с нормой весов задолго до того, как метрика Accuracy фиксирует генерализацию. В нижней строке (без регуляризации, $WD=0.0$) гроккинг не наступает, сеть уходит в вечную меморизацию, и метод MLE корректно показывает не спадающую размерность фазового пространства. Это доказывает отсутствие ложноположительных срабатываний.

Как итог наш метод успешно предсказывает гроккинг, фиксируя топологический коллапс размерности аттрактора. Однако как мы видим, норма весов не является типичным наблюдателем, так как она имеет симетрии, из-за чего она не передаёт исстемную размерность аттрактора, но зато довольно чётко передаёт падение размерности и что важно не зависит от валидационной выборки. Также как мы видим на графиках наш предиктор не превосходит по точности предсказания гроккинга предиктор основанный на норме весов и val\_loss, что связано c простотой задачи.

\begin{figure}[H]
    \centering
    % Верхний ряд: Гроккинг (WD=1.0)
    \begin{subfigure}[b]{0.32\textwidth}
        \includegraphics[width=\textwidth]{images/mod_wd1_acc.png}
        \caption{WD=1: EDM \& Accuracy}
    \end{subfigure}
    \hfill
    \begin{subfigure}[b]{0.32\textwidth}
        \includegraphics[width=\textwidth]{images/mod_wd1_loss.png}
        \caption{WD=1: EDM \& Loss}
    \end{subfigure}
    \hfill
    \begin{subfigure}[b]{0.32\textwidth}
        \includegraphics[width=\textwidth]{images/mod_wd1_norm.png}
        \caption{WD=1: EDM \& Norm}
    \end{subfigure}

    \vspace{0.3cm}
    % Нижний ряд: Нет гроккинга (WD=0.0)
    \begin{subfigure}[b]{0.32\textwidth}
        \includegraphics[width=\textwidth]{images/mod_wd0_acc.png}
        \caption{WD=0: EDM \& Accuracy}
    \end{subfigure}
    \hfill
    \begin{subfigure}[b]{0.32\textwidth}
        \includegraphics[width=\textwidth]{images/mod_wd0_loss.png}
        \caption{WD=0: EDM \& Loss}
    \end{subfigure}
    \hfill
    \begin{subfigure}[b]{0.32\textwidth}
        \includegraphics[width=\textwidth]{images/mod_wd0_norm.png}
        \caption{WD=0: EDM \& Norm}
    \end{subfigure}
    \caption{Этап 2: Анализ гроккинга на модульном сложении при наличии стохастического шума. Сравнение динамики размерности аттрактора построенного по ряду weight\_norm $E$ относительно Accuracy, Loss и weight\_norm в режимах с генерализацией (верхний ряд) и без нее (нижний ряд).}
    \label{fig:grokking_flat}
\end{figure}

\textbf{Этап 3: Неабелевы группы и сложные алгебраические представления ($S_5$).}
Для верификации универсальности метода задача была усложнена до композиции перестановок неабелевой группы $S_5$. Согласно теории представлений, данная математическая структура требует более высоких размерностей для вложения в весовые матрицы трансформера, чем коммутативное сложение.

Как видно на Рис. \ref{fig:grokking_s5}, предложенный метод успешно справился с этой задачей. В фазе хаотичного поиска размерность $E$ держится на высоком уровне. При формировании сетью обобщающего алгебраического правила (гроккинг), размерность предсказуемо коллапсирует. Эксперимент без Weight Decay (нижний ряд) вновь подтверждает, что размерность не падает в случае когда нет гроккинга.

\begin{figure}[H]
    \centering
    % Верхний ряд: Гроккинг S5 (WD=1.0)
    \begin{subfigure}[b]{0.32\textwidth}
        \includegraphics[width=\textwidth]{images/s5_wd1_acc.png}
        \caption{S5 WD=1: Размерность \& Accuracy}
    \end{subfigure}
    \hfill
    \begin{subfigure}[b]{0.32\textwidth}
        \includegraphics[width=\textwidth]{images/s5_wd1_loss.png}
        \caption{S5 WD=1: Размерность \& Loss}
    \end{subfigure}
    \hfill
    \begin{subfigure}[b]{0.32\textwidth}
        \includegraphics[width=\textwidth]{images/s5_wd1_norm.png}
        \caption{S5 WD=1: Размерность \& Norm}
    \end{subfigure}

    \vspace{0.3cm}
    % Нижний ряд: Нет гроккинга S5 (WD=0.0)
    \begin{subfigure}[b]{0.32\textwidth}
        \includegraphics[width=\textwidth]{images/s5_wd0_acc.png}
        \caption{S5 WD=0: Размерность \& Accuracy}
    \end{subfigure}
    \hfill
    \begin{subfigure}[b]{0.32\textwidth}
        \includegraphics[width=\textwidth]{images/s5_wd0_loss.png}
        \caption{S5 WD=0: Размерность \& Loss}
    \end{subfigure}
    \hfill
    \begin{subfigure}[b]{0.32\textwidth}
        \includegraphics[width=\textwidth]{images/s5_wd0_norm.png}
        \caption{S5 WD=0: Размерность \& Norm}
    \end{subfigure}
    \caption{Этап 3: Анализ композиции в группе $S_5$. Коллапс размерности аттрактора построенного по ряду weight\_norm $E$ надежно фиксирует структурную перестройку сети даже при более сложной многомерной структуре скрытых представлений неабелевой группы.}
    \label{fig:grokking_s5}
\end{figure}

Как видно на Рис. \ref{fig:grokking_val_loss_s5}, размерность $E$ восстановленная по валидационной функции потерь также позволяет фиксировать гроккинг, что доказывает, что данный эффект не зависит от выбранной метрики и не является субъективным. Так мы видим ято взяв типичную метрику, мы можем смотреть не только на колапс её размерности, но и на её абсолютные значения, что может быть полезно для понимания количества активных компонентов в модели. В случае $S_5$ мы видим, что размерность коллапсирует до примерно 4, что может быть связано с тем, что для представления $S_5$ в трансформере требуется как минимум 4 компонента. При этом в случае $WD=0$ размерность коллапсирует не падает, а даже растёт, что показывает переобучение модели и уход в вечную меморизацию.


\begin{figure}[H]
    \centering
    \begin{subfigure}[b]{0.49\textwidth}
        \includegraphics[width=\textwidth]{images/s5_wd1_val_loss.pdf}
        \caption{S5 WD=1: Размерность по validation loss}
    \end{subfigure}
    \hfill
    \begin{subfigure}[b]{0.49\textwidth}
        \includegraphics[width=\textwidth]{images/s5_wd0_val_loss.pdf}
        \caption{S5 WD=0: Размерность по validation loss}
    \end{subfigure}

    \vspace{0.3cm}
    \caption{Этап 3: Анализ композиции в группе $S_5$. Размерность $E$ востановленная по validation loss позволяет понимать количество активных компонентов.}
    \label{fig:grokking_val_loss_s5}
\end{figure}

\subsection{Теоретическая интерпретация результатов}\label{sec:exp_theory}

Анализ полученных графиков позволяет сделать следующие теоретические выводы:

\begin{enumerate}
    \item \textbf{Природа гроккинга:} Зафиксированный коллапс размерности (Рис. \ref{fig:grokking_main}b) доказывает, что гроккинг — это не просто уменьшение ошибки, а качественная перестройка динамической системы: переход с высокоэнтропийного аттрактора "запоминания" на низкоразмерный аттрактор "алгебраического правила".
    \item \textbf{Метрики оптимизатора:} Более сложные экспериментв показали, что мы можем диагностировать гроккинг не использую валидационную выборку. Это показывает, что данное явления не является субъективным и не зависит напрямую от выбранной валидационной выборки.
    \item \textbf{Размерность при $wd=0$:} Как мы видим размерность при типичном наблюдателе (validation loss) при $wd=0$ не падает, а даже растёт. Что во первых подтверждает, что при переобучение модель просто заучивает выборку используя большое количество активных компонент, а во вторых показывает важность выбора типичного наблюдателя для оценки алгебраической сложности выученного алгоритма. В случае не типичного наблюдателя (норма весов) размерность на протяжении всего обучения остаётся низкой, так как модель просто находит простое решение и в силу устройства кросс энтропии просто начинает наращивать норму весов из-за чего размерность не растёт.
\end{enumerate}

\section{Обсуждение результатов (Discussion)}\label{sec:disc}

Результаты проведенных вычислительных экспериментов подтверждают фундаментальную гипотезу работы: процесс оптимизации глубоких нейронных сетей может быть исследован с точки зрения теории нелинейных динамических систем и отслеживая размерность аттрактора можно предсказывать и качественно описывать гроккинг. Однако перенос методов фазовой реконструкции в ML выявил ряд специфических ограничений и физических эффектов, требующих отдельного осмысления.

\textbf{1. Проблема вырожденного наблюдателя и роль стохастичности.}
В ходе экспериментов с полнопакетным градиентным спуском (Full-Batch GD) мы столкнулись с сложностью оценки размерности аттрактора: все используемые нами 1-D скалярные метрики вырождались в константу. С точки зрения нелинейной динамики, идеальная оптимизация без шума схлопывает траекторию в прямолинейный луч, лишая методы (такие как MLE) необходимой дисперсии для оценки кривизны ландшафта. Таким образом при анализе размерности стоит следить за тем, чтобы ряд не выраждался. Это можно делать за счёт прореживания ряда, но это может негативно сказаться на точности оценки. Таким образом этот метод работает лучше всего в системах где ряды не выраждаются, то есть в системах со стохастикой. Высокочастотный шум SGD заставляет оптимизатор исследовать локальную геометрию аттрактора. Кроме того, эксперименты показали, что норма весов ($||\vw||_2$) является гораздо более информативным и стабильным «наблюдателем» скрытой динамики, чем склонный к вырождению обучающий лосс.

\textbf{3. Зависимость топологии от выбора наблюдателя (Observability).}
Одним из важнейших эмпирических результатов нашей работы стало выявление зависимости реконструированной размерности $E$ от выбора скалярной метрики (функции наблюдения). В теории динамических систем теорема Такенса требует, чтобы функция наблюдения была «типичной» (generic) и не имела вырожденных симметрий.

Наши эксперименты показали, что норма весов ($||\vw||_2$) является симметричным наблюдателем. Динамика нормы определяется преимущественно балансом двух глобальных сил: градиентом ошибки (увеличивающим веса) и $L_2$-регуляризацией (сжимающей веса). Из-за этого эффективная размерность ряда $||\vw||_2$ всегда коллапсирует до низких значений ($E \approx 2\dots 4$), независимо от сложности решаемой задачи, так как эта метрика «слепа» к сложной угловой структуре представлений.

В противовес этому, функция потерь ($Ls_{val}$ или зашумленный $Ls_{train}$) напрямую зависит от логитов и чувствительна к внутренней алгебраической структуре сети. Работы по прямому спектральному анализу весов \cite{nanda2023progress} показывают, что для модульного сложения сеть формирует 2D-окружности (Фурье-базисы). Оценка MLE по ряду лосса подтверждает это топологически: коллапс траектории останавливается на $E \approx 2$. В свою очередь, для неабелевой группы $S_5$, неприводимые представления которой требуют пространств большей размерности \cite{chughtai2023toy}, коллапс размерности, измеряемый именно по ряду лосса, предсказуемо останавливается на более высоком уровне $E \approx 4\dots 6$. Таким образом, мы доказали, что для «слепой» оценки алгебраической сложности выученного алгоритма необходимо использовать метрики, напрямую связанные с распределением выходов сети, или более аккуратно агрегировать статистики параметров. Также стоит отметить, что лучше использовать именно $Ls_{val}$, так как $Ls_{train}$ быстро выраждается.

\textbf{4. Независимость от валидационной выборки и количество активных компонент.}
С практической точки зрения, мониторинг валидационной функции потерь ($Ls_{val}$) остается более быстрым и дешевым предиктором наступления гроккинга, по крайней мере в простых проведённых экспериментах. Однако метрики являются лишь симптомами улучшения обобщающей способности. В свою очередь, обвал эффективной размерности $E$ предоставляет строгое математическое доказательство структурного фазового перехода — процесса, при котором сеть отсекает шумовые степени свободы и сжимает траекторию в компактное подмногообразие. Также для нашего метода в принципе не нужна валидационная выборка. Это наблюдение открывает путь к принципиально новым методам прунинга (network pruning) и сжатия моделей (например, выбора ранга для LoRA), основанным на истинном количестве степеней свободы обученной сети, вычисленном динамически.

\section{Заключение (Conclusion)}\label{sec:concl}

В данной работе был предложен и верифицирован новый статистический метод для ранней диагностики гроккинга, основанный на топологической реконструкции фазового пространства нейронной сети. Относясь к процессу градиентного спуска как к вынужденной динамической системе, мы пришли к следующим основным результатам:

\begin{itemize}
    \item \textbf{Гроккинг как топологический переход:} Эмпирически доказано, что феномен отложенной генерализации представляет собой структурный фазовый переход системы к низкоразмерному состоянию.

    \item \textbf{Раннее предсказание генерализации:} Использование метода максимального правдоподобия (MLE Intrinsic Dimension) позволило зафиксировать коллапс эффективной размерности аттрактора $E$ задолго до фактического скачка валидационной точности, что утверждает $E$ в качестве предиктора генерализации.

    \item \textbf{Независимость от валидационной выборки:} Показано, что топологический коллапс может быть зафиксирован по ряду метрик, не зависящих от валидационной выборки (например, норма весов), что открывает путь к мониторингу генерализации без необходимости выделения валидационного сета.

    \item \textbf{Универсальность метода:} Метод успешно применялся как к простым коммутативным задачам (модульное сложение), так и к более сложным алгебраическим структурам (композиция в $S_5$), демонстрируя универсальность подхода. Также стоит отметить что данный метод очень просто перенести на новую архитектуру или новую задачу, так как он не требует никаких специфических знаний о внутренней структуре модели и может работать исключительно по скалярам, что делает его очень удобным для практического использования.

    \item \textbf{Алгебраическая сложность алгоритма:} Наблюдение, что при использовании типичной метрики коллапс размерности останавливается на уровнях, соответствующих теоретической сложности задачи (например, $E \approx 2$ для модульного сложения и $E \approx 4\dots 6$ для $S_5$), указывает на то, что размерность аттрактора может служить индикатором алгебраической сложности выученного алгоритма.
\end{itemize}

Главный итог исследования заключается в том, что одномерные скалярные логи (такие как функция потерь, норма весов и другие) позволяют диагностировать скрфтые фазовые переходы нейронной сети во время обучения. Предложенный подход позволяет оценивать истинную алгебраическую сложность выученных алгоритмов исключительно по скалярным рядам, открывая путь к созданию вычислительно дешевых методов мониторинга и структурной оптимизации моделей без необходимости анализа тяжелых весовых матриц или Гессиана.


%%%%%%%%%%%%%%%%%%%%%%%%%%%%%%%%%%%%%%%%%%%%%%%%%%%%%%%%%%%%

\bibliographystyle{unsrtnat}
\bibliography{references}

%%%%%%%%%%%%%%%%%%%%%%%%%%%%%%%%%%%%%%%%%%%%%%%%%%%%%%%%%%%%

% \newpage
% \appendix
% \section{Appendix / supplemental material}\label{app}

% \subsection{Additional experiments / Proofs of Theorems}\label{app:exp}

% TODO

\end{document}